% ============================================================
% Question Bank: ch01-05 Graphing Proportional Relationships
% Topic Title: Graphing Proportional Relationships
% CCSS: 7.RP.A.2d
% Questions: 27 (15 MC + 6 SA + 6 Graphical)
% ============================================================

% --- Multiple Choice Questions (15) ---

\begin{practiceQuestion}{1-5-q01}{mc}
\begin{questionText}
A graph of a proportional relationship passes through the origin and the point $(1, 4)$. What does the point $(1, 4)$ represent?
\end{questionText}
\choiceA{The total for $4$ units}
\choiceB{The unit rate is $4$}
\choiceC{The $y$-intercept is $4$}
\choiceD{The slope is $\frac{1}{4}$}
\correctAnswer{B}
\explanation{On a proportional graph, the point $(1, r)$ gives the unit rate. Here $r = 4$, so the unit rate is $4$ --- meaning $y$ increases by $4$ for every $1$ unit of $x$.}
\end{practiceQuestion}

\begin{practiceQuestion}{1-5-q02}{mc}
\begin{questionText}
A proportional graph passes through the origin and $(5, 20)$. What is the $y$-value when $x = 1$?
\end{questionText}
\choiceA{$1$}
\choiceB{$4$}
\choiceC{$5$}
\choiceD{$20$}
\correctAnswer{B}
\explanation{The unit rate is $k = \frac{20}{5} = 4$. The point $(1, 4)$ is on the graph, so $y = 4$ when $x = 1$.}
\end{practiceQuestion}

\begin{practiceQuestion}{1-5-q03}{mc}
\begin{questionText}
Which statement correctly describes the point $(0, 0)$ on a proportional graph that shows the cost of apples versus weight?
\end{questionText}
\choiceA{It means apples are free.}
\choiceB{It means buying $0$ pounds costs $\$0$.}
\choiceC{It means the graph is curved.}
\choiceD{It means the cost does not change.}
\correctAnswer{B}
\explanation{The origin $(0, 0)$ on a proportional graph means that when the independent variable is $0$, the dependent variable is also $0$. Zero pounds of apples cost zero dollars.}
\end{practiceQuestion}

\begin{practiceQuestion}{1-5-q04}{mc}
\begin{questionText}
Line $M$ passes through $(0, 0)$ and $(3, 15)$. Line $N$ passes through $(0, 0)$ and $(3, 9)$. Which line is steeper?
\end{questionText}
\choiceA{Line $N$, because $9$ is closer to $3$}
\choiceB{Line $M$, because it has a greater unit rate}
\choiceC{They have the same steepness}
\choiceD{Line $N$, because its slope is $3$}
\correctAnswer{B}
\explanation{Line $M$: $k = \frac{15}{3} = 5$. Line $N$: $k = \frac{9}{3} = 3$. A greater unit rate means a steeper line. Since $5 > 3$, Line $M$ is steeper.}
\end{practiceQuestion}

\begin{practiceQuestion}{1-5-q05}{mc}
\begin{questionText}
The graph of $y = 2.5x$ passes through the origin. Which point is also on this graph?
\end{questionText}
\choiceA{$(2, 7.5)$}
\choiceB{$(4, 10)$}
\choiceC{$(3, 6)$}
\choiceD{$(5, 10)$}
\correctAnswer{B}
\explanation{Substitute each $x$ into $y = 2.5x$: at $x = 4$, $y = 2.5 \times 4 = 10$. So $(4, 10)$ is on the graph. The other points do not satisfy the equation.}
\end{practiceQuestion}

\begin{practiceQuestion}{1-5-q06}{mc}
\begin{questionText}
A runner's distance is proportional to time. The graph passes through $(2, 14)$. What is the runner's speed?
\end{questionText}
\choiceA{$2$ miles per hour}
\choiceB{$7$ miles per hour}
\choiceC{$12$ miles per hour}
\choiceD{$14$ miles per hour}
\correctAnswer{B}
\explanation{The unit rate is $k = \frac{14}{2} = 7$ miles per hour. On the graph, the point $(1, 7)$ gives the speed directly.}
\end{practiceQuestion}

\begin{practiceQuestion}{1-5-q07}{mc}
\begin{questionText}
Two students each graph a proportional relationship. Student A's line passes through $(2, 6)$ and Student B's line passes through $(4, 8)$. Both pass through the origin. Which of the following is true?
\end{questionText}
\choiceA{Student A has a greater unit rate and a steeper line.}
\choiceB{Student B has a greater unit rate and a steeper line.}
\choiceC{Both have the same unit rate.}
\choiceD{Student B has a unit rate of $4$.}
\correctAnswer{A}
\explanation{Student A: $k = \frac{6}{2} = 3$. Student B: $k = \frac{8}{4} = 2$. Since $3 > 2$, Student A has the greater unit rate and a steeper line.}
\end{practiceQuestion}

\begin{practiceQuestion}{1-5-q08}{mc}
\begin{questionText}
A proportional graph shows gallons of water ($y$) versus hours ($x$). The line passes through $(3, 10.5)$. How many gallons flow in $8$ hours?
\end{questionText}
\choiceA{$21$}
\choiceB{$24$}
\choiceC{$28$}
\choiceD{$31.5$}
\correctAnswer{C}
\explanation{Find $k = \frac{10.5}{3} = 3.5$ gallons per hour. For $8$ hours: $y = 3.5 \times 8 = 28$ gallons.}
\end{practiceQuestion}

\begin{practiceQuestion}{1-5-q09}{mc}
\begin{questionText}
Which description best fits a graph that is a straight line passing through the origin with a negative slope?
\end{questionText}
\choiceA{A proportional relationship where $y$ increases as $x$ increases}
\choiceB{A proportional relationship where $y$ decreases as $x$ increases}
\choiceC{Not a proportional relationship because the slope is negative}
\choiceD{Not a proportional relationship because it passes through $(0, 0)$}
\correctAnswer{B}
\explanation{A proportional relationship can have a negative constant of proportionality. If $k < 0$, the line still passes through the origin but slopes downward --- $y$ decreases as $x$ increases.}
\end{practiceQuestion}

\begin{practiceQuestion}{1-5-q10}{mc}
\begin{questionText}
The equation $c = 6.50t$ represents the cost of parking for $t$ hours. What does the steepness of the graph tell you?
\end{questionText}
\choiceA{The total cost}
\choiceB{The number of hours parked}
\choiceC{The rate of $\$6.50$ per hour}
\choiceD{The starting cost}
\correctAnswer{C}
\explanation{The steepness (slope) of a proportional graph equals the unit rate $k$. Here $k = 6.50$, so the slope tells you the parking costs $\$6.50$ for each hour.}
\end{practiceQuestion}

\begin{practiceQuestion}{1-5-q11}{mc}
\begin{questionText}
A proportional graph of distance (miles) versus time (hours) passes through $(1, 45)$. A second graph passes through $(1, 60)$. Which driver goes farther in $5$ hours?
\end{questionText}
\choiceA{The first driver, who goes $225$ miles}
\choiceB{The second driver, who goes $300$ miles}
\choiceC{They go the same distance}
\choiceD{The first driver, who goes $250$ miles}
\correctAnswer{B}
\explanation{The point $(1, r)$ gives the unit rate. First driver: $45$ mph, so $5 \times 45 = 225$ miles. Second driver: $60$ mph, so $5 \times 60 = 300$ miles. The second driver goes farther.}
\end{practiceQuestion}

\begin{practiceQuestion}{1-5-q12}{mc}
\begin{questionText}
On a proportional graph, the point $(6, 9)$ is plotted. Which additional point must also be on the graph?
\end{questionText}
\choiceA{$(3, 6)$}
\choiceB{$(0, 0)$}
\choiceC{$(9, 6)$}
\choiceD{$(1, 3)$}
\correctAnswer{B}
\explanation{Every proportional graph passes through the origin $(0, 0)$. This is a defining feature --- zero of one quantity gives zero of the other.}
\end{practiceQuestion}

\begin{practiceQuestion}{1-5-q13}{mc}
\begin{questionText}
A graph shows the cost of buying notebooks. The line passes through $(0, 0)$ and $(4, 14)$. What is the cost of $10$ notebooks?
\end{questionText}
\choiceA{$\$20$}
\choiceB{$\$28$}
\choiceC{$\$35$}
\choiceD{$\$40$}
\correctAnswer{C}
\explanation{Find the unit rate: $k = \frac{14}{4} = 3.50$ per notebook. For $10$ notebooks: $3.50 \times 10 = 35$ dollars.}
\end{practiceQuestion}

\begin{practiceQuestion}{1-5-q14}{mc}
\begin{questionText}
A student says, "The graph of $y = 3x + 1$ is proportional because it is a straight line." What is the student's error?
\end{questionText}
\choiceA{The graph is not a straight line.}
\choiceB{A straight line is only proportional if it passes through the origin, and $y = 3x + 1$ passes through $(0, 1)$.}
\choiceC{The value $3$ is too large to be a slope.}
\choiceD{The graph passes through $(0, 0)$ so the student is correct.}
\correctAnswer{B}
\explanation{A proportional relationship requires a straight line through $(0, 0)$. The equation $y = 3x + 1$ has a $y$-intercept of $1$, so it does not pass through the origin and is not proportional.}
\end{practiceQuestion}

\begin{practiceQuestion}{1-5-q15}{mc}
\begin{questionText}
Two proportional lines are graphed. Line $P$ has a unit rate of $\frac{3}{2}$. Line $Q$ has a unit rate of $\frac{5}{4}$. Which statement is true?
\end{questionText}
\choiceA{Line $Q$ is steeper because $5 > 3$.}
\choiceB{Line $P$ is steeper because $\frac{3}{2} > \frac{5}{4}$.}
\choiceC{They have the same steepness because $3 + 2 = 5 = 5$ and $4 + 1 = 5$.}
\choiceD{Line $Q$ is steeper because $\frac{5}{4} > 1$.}
\correctAnswer{B}
\explanation{Compare: $\frac{3}{2} = 1.5$ and $\frac{5}{4} = 1.25$. Since $1.5 > 1.25$, Line $P$ has the greater unit rate and is steeper.}
\end{practiceQuestion}

% --- Short Answer Questions (6) ---

\begin{practiceQuestion}{1-5-q16}{sa}
\begin{questionText}
A proportional graph passes through $(4, 22)$. What is the unit rate?
\end{questionText}
\correctAnswer{$5.5$}
\explanation{The unit rate is $k = \frac{y}{x} = \frac{22}{4} = 5.5$. The point $(1, 5.5)$ is also on the graph.}
\end{practiceQuestion}

\begin{practiceQuestion}{1-5-q17}{sa}
\begin{questionText}
A proportional graph of cost versus items passes through $(1, 3.75)$. How much do $12$ items cost?
\end{questionText}
\correctAnswer{$\$45$}
\explanation{The unit rate from $(1, 3.75)$ is $\$3.75$ per item. For $12$ items: $3.75 \times 12 = 45$ dollars.}
\end{practiceQuestion}

\begin{practiceQuestion}{1-5-q18}{sa}
\begin{questionText}
Two proportional lines pass through the origin. Line $A$ goes through $(3, 12)$ and Line $B$ goes through $(5, 15)$. What is the difference between their unit rates?
\end{questionText}
\correctAnswer{$1$}
\explanation{Line $A$: $k = \frac{12}{3} = 4$. Line $B$: $k = \frac{15}{5} = 3$. Difference: $4 - 3 = 1$.}
\end{practiceQuestion}

\begin{practiceQuestion}{1-5-q19}{sa}
\begin{questionText}
A proportional graph shows that $6$ pounds of strawberries cost $\$16.50$. What is the cost of $1$ pound?
\end{questionText}
\correctAnswer{$\$2.75$}
\explanation{The unit rate is $k = \frac{16.50}{6} = 2.75$ dollars per pound. This is the $y$-coordinate at $x = 1$ on the graph.}
\end{practiceQuestion}

\begin{practiceQuestion}{1-5-q20}{sa}
\begin{questionText}
Which is steeper: a line through $(0, 0)$ and $(2, 9)$, or a line through $(0, 0)$ and $(3, 12)$? State the unit rate of the steeper line.
\end{questionText}
\correctAnswer{$4.5$}
\explanation{First line: $k = \frac{9}{2} = 4.5$. Second line: $k = \frac{12}{3} = 4$. Since $4.5 > 4$, the first line is steeper with a unit rate of $4.5$.}
\end{practiceQuestion}

\begin{practiceQuestion}{1-5-q21}{sa}
\begin{questionText}
A proportional graph passes through $(8, 6)$. What is the value of $y$ when $x = 20$?
\end{questionText}
\correctAnswer{$15$}
\explanation{Find $k = \frac{6}{8} = 0.75$. Then $y = 0.75 \times 20 = 15$.}
\end{practiceQuestion}

% --- Graphical Multiple Choice Questions (3) ---

\begin{practiceQuestion}{1-5-q22}{gmc}
\begin{questionText}
The graph below shows the distance a train travels over time.

\begin{center}
\begin{tikzpicture}[scale=0.5]
  \draw[step=1,gray!20,thin] (0,0) grid (8,10);
  \draw[thick,->] (0,0) -- (8.5,0) node[right]{\small Hours};
  \draw[thick,->] (0,0) -- (0,10.5) node[above]{\small Distance (mi)};
  \foreach \x in {1,...,8} \draw (\x,0.1) -- (\x,-0.1) node[below]{\tiny \x};
  \foreach \y/\lab in {1/50,2/100,3/150,4/200,5/250,6/300,7/350,8/400,9/450,10/500}
    \draw (0.1,\y) -- (-0.1,\y) node[left]{\tiny \lab};
  \draw[thick,funBlue] (0,0) -- (8,6.4);
  \fill[funBlue] (0,0) circle (3pt);
  \fill[funBlue] (1,0.8) circle (3pt);
  \fill[funBlue] (5,4) circle (3pt);
\end{tikzpicture}
\end{center}

What is the train's speed in miles per hour?
\end{questionText}
\choiceA{$40$ mph}
\choiceB{$50$ mph}
\choiceC{$80$ mph}
\choiceD{$100$ mph}
\correctAnswer{A}
\explanation{From the point $(1, 0.8)$ on the graph, and since each grid unit on the $y$-axis equals $50$ miles, the $y$-value at $x = 1$ is $0.8 \times 50 = 40$ miles. The unit rate from the point $(5, 200)$ confirms: $200 \div 5 = 40$ mph.}
\end{practiceQuestion}

\begin{practiceQuestion}{1-5-q23}{gmc}
\begin{questionText}
Two proportional lines are graphed below. Line $A$ (blue) passes through $(4, 6)$. Line $B$ (orange) passes through $(4, 10)$.

\begin{center}
\begin{tikzpicture}[scale=0.5]
  \draw[step=1,gray!20,thin] (0,0) grid (8,12);
  \draw[thick,->] (0,0) -- (8.5,0) node[right]{$x$};
  \draw[thick,->] (0,0) -- (0,12.5) node[above]{$y$};
  \foreach \x in {1,...,8} \draw (\x,0.1) -- (\x,-0.1) node[below]{\tiny \x};
  \foreach \y in {2,4,...,12} \draw (0.1,\y) -- (-0.1,\y) node[left]{\tiny \y};
  \draw[thick,funBlue] (0,0) -- (8,12);
  \draw[thick,funOrange] (0,0) -- (5,12.5);
  \fill[funBlue] (4,6) circle (3pt) node[right=2pt]{\tiny $A$};
  \fill[funOrange] (4,10) circle (3pt) node[left=2pt]{\tiny $B$};
\end{tikzpicture}
\end{center}

What is the unit rate of Line $B$?
\end{questionText}
\choiceA{$1.5$}
\choiceB{$2$}
\choiceC{$2.5$}
\choiceD{$4$}
\correctAnswer{C}
\explanation{Line $B$ passes through $(4, 10)$ and the origin: $k = \frac{10}{4} = 2.5$. The unit rate is $2.5$.}
\end{practiceQuestion}

\begin{practiceQuestion}{1-5-q24}{gmc}
\begin{questionText}
The graph below represents a proportional relationship between cups of flour and number of cakes baked.

\begin{center}
\begin{tikzpicture}[scale=0.5]
  \draw[step=1,gray!20,thin] (0,0) grid (10,8);
  \draw[thick,->] (0,0) -- (10.5,0) node[right]{\small Cakes};
  \draw[thick,->] (0,0) -- (0,8.5) node[above]{\small Cups of flour};
  \foreach \x in {1,...,10} \draw (\x,0.1) -- (\x,-0.1) node[below]{\tiny \x};
  \foreach \y in {1,...,8} \draw (0.1,\y) -- (-0.1,\y) node[left]{\tiny \y};
  \draw[thick,funGreen] (0,0) -- (10,7.5);
  \fill[funGreen] (0,0) circle (3pt);
  \fill[funGreen] (4,3) circle (3pt);
  \fill[funGreen] (8,6) circle (3pt);
\end{tikzpicture}
\end{center}

How many cups of flour are needed for $12$ cakes?
\end{questionText}
\choiceA{$6$}
\choiceB{$8$}
\choiceC{$9$}
\choiceD{$12$}
\correctAnswer{C}
\explanation{From the point $(4, 3)$: $k = \frac{3}{4} = 0.75$ cups per cake. For $12$ cakes: $0.75 \times 12 = 9$ cups.}
\end{practiceQuestion}

% --- Graphical Short Answer Questions (3) ---

\begin{practiceQuestion}{1-5-q25}{gsa}
\begin{questionText}
The graph below shows two proportional lines.

\begin{center}
\begin{tikzpicture}[scale=0.5]
  \draw[step=1,gray!20,thin] (0,0) grid (8,10);
  \draw[thick,->] (0,0) -- (8.5,0) node[right]{$x$};
  \draw[thick,->] (0,0) -- (0,10.5) node[above]{$y$};
  \foreach \x in {1,...,8} \draw (\x,0.1) -- (\x,-0.1) node[below]{\tiny \x};
  \foreach \y in {1,...,10} \draw (0.1,\y) -- (-0.1,\y) node[left]{\tiny \y};
  \draw[thick,funBlue] (0,0) -- (8,8);
  \draw[thick,funRed] (0,0) -- (5,10);
  \fill[funBlue] (4,4) circle (3pt) node[below right]{\tiny $A(4,4)$};
  \fill[funRed] (5,10) circle (3pt) node[left]{\tiny $B(5,10)$};
\end{tikzpicture}
\end{center}

What is the difference in the unit rates of Line $B$ and Line $A$?
\end{questionText}
\correctAnswer{$1$}
\explanation{Line $A$: $k = \frac{4}{4} = 1$. Line $B$: $k = \frac{10}{5} = 2$. Difference: $2 - 1 = 1$.}
\end{practiceQuestion}

\begin{practiceQuestion}{1-5-q26}{gsa}
\begin{questionText}
The coordinate plane below shows a proportional relationship between time (seconds) and distance (meters).

\begin{center}
\begin{tikzpicture}[scale=0.55]
  \draw[step=1,gray!20,thin] (0,0) grid (8,10);
  \draw[thick,->] (0,0) -- (8.5,0) node[right]{\small Seconds};
  \draw[thick,->] (0,0) -- (0,10.5) node[above]{\small Meters};
  \foreach \x in {1,...,8} \draw (\x,0.1) -- (\x,-0.1) node[below]{\tiny \x};
  \foreach \y in {1,...,10} \draw (0.1,\y) -- (-0.1,\y) node[left]{\tiny \y};
  \draw[thick,funOrange] (0,0) -- (8,7);
  \fill[funOrange] (0,0) circle (3pt);
  \fill[funOrange] (4,3.5) circle (3pt);
  \fill[funOrange] (8,7) circle (3pt);
\end{tikzpicture}
\end{center}

What is the speed in meters per second? How far does the object travel in $12$ seconds?
\end{questionText}
\correctAnswer{$10.5$ meters}
\explanation{From $(8, 7)$: $k = \frac{7}{8} = 0.875$ m/s. For $12$ seconds: $0.875 \times 12 = 10.5$ meters.}
\end{practiceQuestion}

\begin{practiceQuestion}{1-5-q27}{gsa}
\begin{questionText}
The graph below shows two workers' earnings over time. Both pass through the origin.

\begin{center}
\begin{tikzpicture}[scale=0.5]
  \draw[step=1,gray!20,thin] (0,0) grid (8,10);
  \draw[thick,->] (0,0) -- (8.5,0) node[right]{\small Hours};
  \draw[thick,->] (0,0) -- (0,10.5) node[above]{\small Earnings (\$)};
  \foreach \x in {1,...,8} \draw (\x,0.1) -- (\x,-0.1) node[below]{\tiny \x};
  \foreach \y/\lab in {2/20,4/40,6/60,8/80,10/100}
    \draw (0.1,\y) -- (-0.1,\y) node[left]{\tiny \lab};
  \draw[thick,funBlue] (0,0) -- (8,9.6);
  \draw[thick,funRed] (0,0) -- (8,6.4);
  \fill[funBlue] (5,6) circle (3pt) node[above left]{\tiny Worker $A$};
  \fill[funRed] (5,4) circle (3pt) node[below right]{\tiny Worker $B$};
\end{tikzpicture}
\end{center}

Each $y$-axis unit equals $\$10$. How much more does Worker $A$ earn than Worker $B$ in an $8$-hour shift?
\end{questionText}
\correctAnswer{$\$32$}
\explanation{Each $y$-unit is $\$10$, so Worker $A$ at $(5, 6)$ earns $\$60$ in $5$ hours ($\$12$/hr) and Worker $B$ at $(5, 4)$ earns $\$40$ in $5$ hours ($\$8$/hr). In $8$ hours: $A$ earns $12 \times 8 = \$96$ and $B$ earns $8 \times 8 = \$64$. Difference: $96 - 64 = \$32$.}
\end{practiceQuestion}
